\chapter{Qubits, Complex State Vectors, and Observables}

\lectureoverview{A single electron spin is the first genuinely quantum bit. It may be prepared relative to any spatial direction, yet a fixed detector returns only two outcomes. We let the classical magnetic picture fail, then construct the complex vector space, probability rule, and matrix language needed to describe what the experiment actually does.}

\section{One Spin and Two Meanings of Vector}

A classical bit has two possible states and nothing between them. We may call them \(0\) and \(1\), heads and tails, or up and down; the names are labels for two distinct points in a classical state space. A quantum bit also yields two alternatives when measured in a fixed way, but it cannot be described by only two state-space points. The simplest example is the spin of an electron.

There are two different notions of vector in this example, and confusing them makes everything that follows unnecessarily mysterious. First there is an ordinary \emph{spatial vector}: an arrow in three-dimensional space with a magnitude and a direction. The electron has a magnetic moment that behaves, for the present purpose, like a tiny magnet. Its magnitude is fixed, while a spatial direction can be assigned to it in a preparation procedure.

The second object is an abstract vector in a complex vector space. This \emph{state vector} will describe the quantum state. It is not simply the magnetic arrow written in another notation. The relation between spatial directions and state vectors is real and important, but it is not the identity of two objects. We shall earn that relation gradually.

\section{Letting the Classical Picture Fail}

Begin deliberately with a classical story. Put a tiny magnet in a strong external magnetic field. Unless it is already aligned, it precesses about the field. As it radiates energy, the precession damps and the magnet settles into its lowest-energy orientation. By choosing the direction of the applied field, waiting for this relaxation, and then removing the field, we seem able to prepare the magnetic moment in any desired direction.

The same apparatus suggests a method of detection. Turn on a field in a chosen direction and observe the radiation produced while the little magnet settles. A classical magnet aligned with the field should emit no energy. One initially opposed to the field should emit the maximum amount, and an intermediate angle should produce an intermediate amount. In that picture the radiated energy varies continuously with the angle.

\begin{figure}[htbp]
\centering
\begin{coursefitfigure}
\begin{tikzpicture}[>=Latex,thick,scale=0.92]
  \draw[rounded corners] (-0.2,-0.2) rectangle (3.0,3.4);
  \node at (1.4,3.05) {\small preparation};
  \draw[->,very thick] (1.4,0.45) -- (1.4,2.55);
  \node[right] at (1.4,2.35) {\small applied field};
  \fill (1.4,1.35) circle (2.3pt);
  \draw[->,very thick] (1.4,1.35) -- (2.25,1.95);
  \node[below] at (1.4,0.15) {\small choose a direction and wait};

  \draw[->] (3.25,1.6) -- (4.25,1.6);

  \draw[rounded corners] (4.5,-0.2) rectangle (8.8,3.4);
  \node at (6.65,3.05) {\small measurement in a fixed direction};
  \draw[->,very thick] (5.15,0.45) -- (5.15,2.55);
  \fill (5.15,1.35) circle (2.3pt);
  \node[anchor=west] at (5.7,2.0) {no photon};
  \node[anchor=west] at (5.7,1.25) {one photon of fixed energy};
\node[anchor=west] at (5.7,0.5) {no intermediate outcome};
\end{tikzpicture}
\end{coursefitfigure}
\caption{The operational puzzle. A continuously variable preparation is followed by a detector with only two possible outputs.}
\end{figure}

That classical prediction is false for an electron. When the detector field is turned on, only two results occur:
\begin{enumerate}
  \item no photon is emitted;
  \item one photon is emitted with the fixed energy separating the down and up levels.
\end{enumerate}
There is never a photon carrying an arbitrary fraction of that energy. The continuous classical response has collapsed into a binary quantum answer.

The direction used in preparation has not become irrelevant. It controls the \emph{probabilities} of the two answers. Prepare the electron up relative to the detector and the up result occurs with certainty. Prepare it down and the transition occurs with certainty. Prepare it horizontally and a vertical detector returns the two answers equally often:
\begin{equation}
  P(\text{up})=\frac12,
  \qquad
  P(\text{down})=\frac12.
  \label{eq:l02-horizontal-probabilities}
\end{equation}
Thus a continuously variable family of preparations leads to a two-valued measurement whose statistics vary continuously.

\begin{classroomqa}
\lecturetimestamp{00:21:23}
\audiencequestion{Could another small object, perhaps a second electron, supply the magnetic field instead of a macroscopic magnet?}
\lecturerresponse{It could, but then the two electrons form a joint quantum system and can become entangled. That is precisely where the course is going. We postpone it until the state space of one electron is under control.}
\end{classroomqa}

\begin{classroomqa}
\lecturetimestamp{00:22:17}
\audiencequestion{Must the probability distribution be found by repeating the entire preparation and measurement? What happens if we measure the same electron again?}
\lecturerresponse{Yes, the original distribution requires repeated, identically prepared trials. After the first measurement, repeating the same measurement on that electron gives the same result; the measurement has prepared it relative to that detector. Waiting for a radiative transition is itself one possible preparation procedure.}
\end{classroomqa}

\begin{classroomqa}
\lecturetimestamp{00:23:48}
\audiencequestion{Are we imagining a free electron, and is the magnetic field left on continuously?}
\lecturerresponse{We isolate the spin degree of freedom and ignore the electron's position, whether it is freely moving or otherwise constrained. A large electromagnet may be switched on quickly for preparation or detection. The simplified experiment is designed to ask only about spin.}
\end{classroomqa}

\begin{classroomqa}
\lecturetimestamp{00:25:22}
\audiencequestion{If the electron is prepared exactly down and the field is reversed, is a photon emitted with certainty, and where does its energy come from?}
\lecturerresponse{For exact opposition the transition occurs with probability one. For an oblique preparation the probability is smaller. The energy ultimately comes from the work done in establishing or reversing the external magnetic field; the apparatus is part of the energy accounting.}
\end{classroomqa}

\begin{classroomqa}
\lecturetimestamp{00:27:31}
\audiencequestion{When does the photon appear? Is the transition instantaneous?}
\lecturerresponse{The emission time is random, like the decay time of an unstable state. One can specify a characteristic lifetime or half-life, and a stronger field shortens that scale. Energy--time uncertainty enters the finer description, but the present operational facts are the two possible outputs and their measured frequencies.}
\end{classroomqa}

\begin{classroomqa}
\lecturetimestamp{00:29:17}
\audiencequestion{Would a beam containing a million identically prepared electrons simply produce a measurable light intensity?}
\lecturerresponse{Yes. The ensemble intensity reveals the average transition probability, while the emission time of any particular electron remains random. An ensemble displays the distribution that a single trial cannot.}
\end{classroomqa}

This is the point of departure. Two isolated classical points can encode the detector answers, but they cannot encode all the preparations that alter the probabilities of those answers. We need a different state space.

\section{Complex Numbers and the Geometry of State Space}

The required state space is linear: its elements can be added and multiplied by scalars. Its scalars are complex numbers. Write
\begin{equation}
  z=a+ib,
  \qquad
  i^2=-1,
  \qquad
  z^*=a-ib,
  \label{eq:l02-complex-number}
\end{equation}
where \(a\) and \(b\) are real. Complex conjugation produces the positive quantity
\begin{align}
  z^*z
  &=(a-ib)(a+ib) \notag\\
  &=a^2+b^2.
  \label{eq:l02-complex-norm}
\end{align}
This simple cancellation is the algebraic seed of the probability rule.

\begin{figure}[htbp]
\centering
\includegraphics[width=0.74\linewidth]{lecture_02_figure_02.png}
\caption{Complex conjugation and \(z^*z\) are developed on the board before amplitudes are interpreted as probabilities.}
\lectureframe{00:37:10}
\end{figure}

A complex vector space contains objects \(\ket a,\ket b,\ldots\) for which
\begin{equation}
  c\ket a,
  \qquad
  \ket a+\ket b
\end{equation}
are again vectors whenever \(c\) is complex. Column vectors give a concrete realization:
\begin{equation}
  \begin{aligned}
    \ket a&=
    \begin{pmatrix}
      a_1\\a_2\\ \vdots\\a_n
    \end{pmatrix},
    &
    c\ket a&=
    \begin{pmatrix}
      ca_1\\ca_2\\ \vdots\\ca_n
    \end{pmatrix},\\[0.7em]
    \ket a+\ket b&=
    \begin{pmatrix}
      a_1+b_1\\a_2+b_2\\ \vdots\\a_n+b_n
    \end{pmatrix}.
  \end{aligned}
  \label{eq:l02-vector-operations}
\end{equation}
Even the complex numbers themselves form a one-dimensional complex vector space. For an electron spin we shall need two dimensions.

To each column ket we associate a conjugate row bra:
\begin{equation}
  \ket a=
  \begin{pmatrix}a_1\\a_2\end{pmatrix}
  \quad\longleftrightarrow\quad
  \bra a=
  \begin{pmatrix}a_1^*&a_2^*\end{pmatrix}.
  \label{eq:l02-bra-ket-components}
\end{equation}
The inner product is then
\begin{equation}
  \braket{b}{a}=b_1^*a_1+b_2^*a_2,
  \label{eq:l02-inner-product}
\end{equation}
and the squared norm is
\begin{equation}
  \braket{a}{a}=|a_1|^2+|a_2|^2.
  \label{eq:l02-norm}
\end{equation}
It is real and nonnegative because every term has the form \(z^*z\).

\begin{classroomqa}
\lecturetimestamp{00:45:24}
\audiencequestion{Why is there no explicit star on the bra symbol?}
\lecturerresponse{Turning a ket into its bra already means transpose and complex conjugate. The reversal of the bracket is notation for that operation, so writing another star would repeat what is implicit.}
\end{classroomqa}

\section{The Qubit Basis and the Born Rule}

Return now to the detector. Let \(\ket{+}\) denote its up state and \(\ket{-}\) its down state. A general spin state is a linear combination
\begin{equation}
  \ket\psi=a_+\ket{+}+a_-\ket{-}
  \quad\longleftrightarrow\quad
  \begin{pmatrix}a_+\\a_-\end{pmatrix}.
  \label{eq:l02-general-qubit}
\end{equation}
The plus and minus signs inside the kets are labels, not arithmetic operations. Once this ordered basis is chosen,
\begin{equation}
  \ket{+}=
  \begin{pmatrix}1\\0\end{pmatrix},
  \qquad
  \ket{-}=
  \begin{pmatrix}0\\1\end{pmatrix}.
  \label{eq:l02-basis-columns}
\end{equation}
Indeed,
\begin{equation}
  a_+\begin{pmatrix}1\\0\end{pmatrix}
  +a_-\begin{pmatrix}0\\1\end{pmatrix}
  =
  \begin{pmatrix}a_+\\a_-\end{pmatrix}.
  \label{eq:l02-basis-expansion}
\end{equation}

\begin{figure}[htbp]
\centering
\includegraphics[width=0.72\linewidth]{lecture_02_figure_03.png}
\caption{The detector states are fixed as the two standard basis columns.}
\lectureframe{00:58:10}
\end{figure}

The coefficients are probability amplitudes. The probability of each detector outcome is the squared magnitude of the corresponding amplitude:
\begin{equation}
  P(+)=|a_+|^2,
  \qquad
  P(-)=|a_-|^2.
  \label{eq:l02-born-rule}
\end{equation}
Because these are the only outcomes,
\begin{equation}
  \begin{gathered}
    |a_+|^2+|a_-|^2=1,\\
    \text{or equivalently}\qquad \braket{\psi}{\psi}=1.
  \end{gathered}
  \label{eq:l02-normalization}
\end{equation}
A physical state vector is therefore normalized.

\begin{classroomqa}
\lecturetimestamp{00:56:54}
\audiencequestion{Do the plus and minus signs have an algebraic meaning in the components?}
\lecturerresponse{Here they label the two basis states and their coefficients. We could call them \(1\) and \(2\), or up and down. In particular, \(-\ket{+}\) is a valid scalar multiple of \(\ket{+}\); it is not the state \(\ket{-}\).}
\end{classroomqa}

Multiplying a state by a complex number of unit magnitude leaves all its probabilities unchanged:
\begin{equation}
  \ket\psi\longmapsto e^{i\phi}\ket\psi,
  \qquad
  |e^{i\phi}a_\pm|^2=|a_\pm|^2.
  \label{eq:l02-global-phase}
\end{equation}
Thus two vectors differing only by an overall phase represent the same physical state. A \emph{relative} phase is different. Replacing only \(a_-\) by \(-a_-\), for example, cannot generally be removed by multiplying the whole vector by one phase; it can change the results of measurements in another basis.

\begin{classroomqa}
\lecturetimestamp{00:59:54}
\audiencequestion{Are \(\ket{+}\), \(-\ket{+}\), and \(e^{i\phi}\ket{+}\) different states? What about changing the sign of one term in a superposition?}
\lecturerresponse{They are different mathematical vectors but, when the factor multiplies the entire vector and has unit magnitude, they describe the same physical state. A relative sign or phase between the up and down terms is observable in suitable measurements and therefore changes the state.}
\end{classroomqa}

The basis is orthonormal:
\begin{align}
  \braket{+}{+}&=1,
  &\braket{-}{-}&=1,
  \notag\\
  \braket{+}{-}&=0,
  &\braket{-}{+}&=0.
  \label{eq:l02-orthonormality}
\end{align}
For example,
\begin{equation}
  \braket{+}{-}
  =
  \begin{pmatrix}1&0\end{pmatrix}
  \begin{pmatrix}0\\1\end{pmatrix}
  =0.
\end{equation}

\begin{figure}[htbp]
\centering
\includegraphics[width=0.76\linewidth]{lecture_02_figure_04.png}
\caption{The four inner products of the up/down basis are checked directly on the board.}
\lectureframe{01:05:00}
\end{figure}

An equal-weight state provides the first useful correction at the board. The column
\begin{equation}
  \begin{pmatrix}1\\1\end{pmatrix}
\end{equation}
has squared norm \(2\), not \(1\), and therefore is not normalized. Dividing by \(\sqrt2\) repairs it:
\begin{equation}
  \ket{\rightarrow}
  =
  \frac{1}{\sqrt2}
  \begin{pmatrix}1\\1\end{pmatrix}
  =
  \frac{\ket{+}+\ket{-}}{\sqrt2}.
  \label{eq:l02-horizontal-state}
\end{equation}
It gives \(P(+)=P(-)=1/2\) in the vertical detector and corresponds to one horizontal preparation. The correction from \((1,1)^T\) to \(2^{-1/2}(1,1)^T\) is not cosmetic: normalization is what turns squared amplitudes into a complete probability distribution.

After the break it is worth returning to the two meanings of vector. A spatial arrow has three real components. A qubit state vector has two complex components, hence four real numbers before imposing normalization and identifying overall phase. The objects are related, but one is not merely a relabeling of the other. We do not yet need the full geometric map between them.

\begin{classroomqa}
\lecturetimestamp{01:10:56}
\audiencequestion{How do we know the dimension of the state space for a different quantum system?}
\lecturerresponse{Experiment tells us how many mutually exclusive outcomes a complete measurement can distinguish. Electron spin has two such outcomes and therefore a two-dimensional complex state space. A spin-one system has three outcomes in the corresponding measurement and requires three components.}
\end{classroomqa}

\section{Classical Observables and Ensemble Averages}

Having described states, turn to measurable quantities. In a classical system with discrete states \(n\), an observable is a real-valued function \(f_n\) on those states. If the state is uncertain and described by probabilities \(p_n\), its ensemble average is
\begin{equation}
  \langle f\rangle_{\mathrm{ens}}
  =\sum_n p_n f_n.
  \label{eq:l02-classical-expectation}
\end{equation}
The notation does not imply that a single measurement must ever equal this number.

For a coin, assign \(f_+=+1\) and \(f_-=-1\). Then
\begin{equation}
  \langle f\rangle_{\mathrm{ens}}
  =p_+-p_-.
  \label{eq:l02-coin-expectation}
\end{equation}
A fair coin has average zero. If \(p_+=3/4\) and \(p_-=1/4\), the average is \(1/2\). Neither \(0\) nor \(1/2\) is an outcome of a single toss; each summarizes many \(+1\) and \(-1\) results.

\begin{classroomqa}
\lecturetimestamp{01:23:30}
\audiencequestion{How can the expected value be zero or one half when neither value can appear in a single trial?}
\lecturerresponse{The phrase expected value is potentially misleading. It means the average of a large ensemble, not the value one should expect to see on the next run. A fair sequence of \(+1\) and \(-1\) results averages to zero, and a three-to-one imbalance averages to one half.}
\end{classroomqa}

An especially useful observable is the indicator for one state:
\begin{equation}
  \chi_k(n)=
  \begin{cases}
    1,&n=k,\\
    0,&n\ne k.
  \end{cases}
\end{equation}
Its average is
\begin{equation}
  \langle\chi_k\rangle_{\mathrm{ens}}=p_k.
\end{equation}
Thus knowledge of the averages of all classical observables includes, in particular, knowledge of the probability distribution.

\section{Matrices and Quantum Observables}

In quantum mechanics observables are represented not by functions on state-space points but by linear operators. In a finite-dimensional state space, those operators are matrices. A matrix \(M\) maps one vector to another:
\begin{equation}
  M\ket a=\ket{a'}.
  \label{eq:l02-matrix-action}
\end{equation}
For two components,
\begin{equation}
  \begin{pmatrix}
    m_{11}&m_{12}\\
    m_{21}&m_{22}
  \end{pmatrix}
  \begin{pmatrix}a_1\\a_2\end{pmatrix}
  =
  \begin{pmatrix}
    m_{11}a_1+m_{12}a_2\\
    m_{21}a_1+m_{22}a_2
  \end{pmatrix}.
  \label{eq:l02-matrix-components}
\end{equation}

\begin{figure}[htbp]
\centering
\includegraphics[width=0.78\linewidth]{lecture_02_figure_05.png}
\caption{Matrix action is expanded row by row into the two components of the transformed vector.}
\lectureframe{01:32:10}
\end{figure}

Real two-dimensional vectors make the geometry visible. The matrices
\begin{equation}
  \begin{pmatrix}2&0\\0&2\end{pmatrix},
  \qquad
  \begin{pmatrix}1&0\\0&2\end{pmatrix},
  \qquad
  \begin{pmatrix}2&0\\0&1\end{pmatrix}
\end{equation}
respectively stretch every direction uniformly, stretch only the \(y\)-component, and stretch only the \(x\)-component. The matrix
\begin{equation}
  \begin{pmatrix}0&1\\1&0\end{pmatrix}
\end{equation}
interchanges the components and reflects the plane across \(x=y\). Meanwhile,
\begin{equation}
  \begin{pmatrix}0&1\\-1&0\end{pmatrix}
  \begin{pmatrix}v_x\\v_y\end{pmatrix}
  =
  \begin{pmatrix}v_y\\-v_x\end{pmatrix}
  \label{eq:l02-quarter-turn}
\end{equation}
rotates vectors through a clockwise quarter-turn in the convention used at the board. A shear supplies another linear example, although its matrix is left as practice. The essential restriction is linearity: not every imaginable transformation of the plane is represented by a matrix.

Quantum observables belong to a special class.

\begin{definition}
A matrix \(M\) is \emph{Hermitian} when
\begin{equation}
  M_{ij}=M_{ji}^*.
  \label{eq:l02-hermitian-condition}
\end{equation}
\end{definition}

The diagonal condition \(M_{ii}=M_{ii}^*\) forces every diagonal entry to be real. Off-diagonal entries occur in conjugate pairs. For example,
\begin{equation}
  M=
  \begin{pmatrix}
    r&4+i\\
    4-i&s
  \end{pmatrix},
  \qquad r,s\in\mathbb R,
  \label{eq:l02-hermitian-example}
\end{equation}
is Hermitian. If all entries happen to be real, Hermitian is the same as symmetric.

\begin{figure}[htbp]
\centering
\includegraphics[width=0.75\linewidth]{lecture_02_figure_06.png}
\caption{The defining conjugation across the diagonal, \(M_{ij}=M_{ji}^*\), is isolated at the end of the lecture.}
\lectureframe{01:42:50}
\end{figure}

Hermitian matrices will represent measurable quantities, including spin along the \(x\), \(y\), and \(z\) directions. They should not be confused with the operators that describe a physical kick or time evolution.

\begin{classroomqa}
\lecturetimestamp{01:45:49}
\audiencequestion{If a matrix kicks the state into another state, must it preserve normalization? Does a Hermitian matrix do that?}
\lecturerresponse{A physical kick must preserve the length of the state vector, but a general Hermitian matrix does not. Length-preserving transformations are unitary. Hermitian operators represent observables; unitary operators represent the reversible transformations that act on states. Their detailed relation comes next.}
\end{classroomqa}

We have reached the minimum language needed for the next step. A qubit is a normalized vector in a two-dimensional complex vector space; amplitudes determine probabilities; and observables are Hermitian matrices. With those pieces in place, spin along arbitrary directions and eventually entangled pairs can be discussed without returning to the classical picture that has already failed.
